\documentclass{beamer}
\usetheme{Madrid}
\usecolortheme{default}
\usepackage{listings}
\usepackage{xcolor}

\title{Proyecto 1 - OLTP - Colecciones}
\subtitle{Bases de Datos a Gran Escala}
\author{Equipo 13: \\ Andres Lopez Gonzalez \\ Carlos Fabian Esteva Gallegos}
\date{}

\lstset{
  language=SQL,
  basicstyle=\tiny\ttfamily,
  keywordstyle=\color{blue}\bfseries,
  commentstyle=\color{green!40!black},
  stringstyle=\color{red},
  breaklines=true,
  showstringspaces=false,
  frame=single,
  numbers=left,
  numberstyle=\tiny,
  stepnumber=1
}

\begin{document}

%---------------------------
\begin{frame}
  \titlepage
\end{frame}

%---------------------------
\begin{frame}{Descripción del Proyecto}
  \begin{itemize}
    \item Encuentro mundial de coleccionistas de tarjetas deportivas.
    \item Se diseña un sistema OLTP para gestionar deportes, equipos, jugadores y tarjetas.
    \item Información clave:
    \begin{itemize}
      \item Deporte: nombre, descripción.
      \item Equipo: nombre, ciudad, fechas de partidos.
      \item Jugador: nombre, género, puesto, fechas, lugar de nacimiento.
      \item Tarjeta: número, estado, fecha venta, medio adquisición.
    \end{itemize}
  \end{itemize}
\end{frame}

%---------------------------
\begin{frame}{Modelo y Reglas}
  \begin{itemize}
    \item Se aplicaron \textbf{constraints} para integridad:
    \begin{itemize}
      \item Validación de nombres, géneros y estados.
      \item Reglas de consistencia en fechas.
      \item Restricciones de adquisición y estado de tarjetas.
    \end{itemize}
    \item Procedimientos almacenados:
    \begin{itemize}
      \item \texttt{sp\_inserta}: inserción de deportes.
      \item \texttt{sp\_borra}: eliminación segura de deportes.
    \end{itemize}
  \end{itemize}
\end{frame}

%---------------------------
\begin{frame}{Consultas y Resultados}
  \begin{itemize}
    \item \texttt{sp\_consultas} implementa 5 consultas parametrizadas:
    \begin{enumerate}
      \item Estado de tarjetas por equipo.
      \item Medio de adquisición por jugador.
      \item Relación deporte-equipo-tarjeta.
      \item Equipos según estado de colección.
      \item Jugadores y puesto dentro de un equipo.
    \end{enumerate}
    \item Resultados exportados a archivos \texttt{.txt} para análisis.
  \end{itemize}
\end{frame}

%---------------------------
\begin{frame}[fragile]{Constraints principales}
\begin{lstlisting}
-- Ejemplo de constraints
ALTER TABLE deporte 
ADD CONSTRAINT chk_nombre_deporte 
CHECK (nombre_deporte ~ '^[A-Za-zÁÉÍÓÚáéíóúÑñ\s]+$');

ALTER TABLE jugador 
ADD CONSTRAINT chk_genero 
CHECK (genero IN ('Masculino','Femenino','Otro'));

ALTER TABLE tarjeta 
ADD CONSTRAINT chk_estado_coleccion 
CHECK (estado_coleccion IN ('Nuevo','Casi Nuevo','Usado',
                            'Daños Visibles','Daño Total'));
\end{lstlisting}
\end{frame}

%---------------------------
\begin{frame}[fragile]{Stored Procedures: Inserción y Borrado}
\begin{lstlisting}
-- Inserción en tabla Deporte
CREATE OR REPLACE FUNCTION sp_inserta_deporte(
    p_nombre_deporte VARCHAR(50),
    p_descripcion TEXT)
RETURNS void AS $$
BEGIN
  INSERT INTO deporte (nombre_deporte, descripcion)
  VALUES (p_nombre_deporte, p_descripcion);
END;
$$ LANGUAGE plpgsql;

-- Borrado seguro en tabla Deporte
CREATE OR REPLACE FUNCTION sp_borra_deporte(p_id_deporte INT)
RETURNS void AS $$
BEGIN
  IF EXISTS (SELECT 1 FROM equipo WHERE id_deporte = p_id_deporte) THEN
    RAISE EXCEPTION 'No se puede eliminar, deporte en uso.';
  END IF;
  DELETE FROM deporte WHERE id_deporte = p_id_deporte;
END;
$$ LANGUAGE plpgsql;
\end{lstlisting}
\end{frame}


%---------------------------
\begin{frame}[fragile]{Stored Procedure: Consultas}
\begin{lstlisting}
-- Procedimiento que ejecuta 5 consultas
CREATE OR REPLACE FUNCTION sp_consultas(
    opcion INT,
    p_nombre_equipo VARCHAR DEFAULT NULL,
    p_nombre_jugador VARCHAR DEFAULT NULL,
    p_nombre_deporte VARCHAR DEFAULT NULL,
    p_estado_coleccion VARCHAR DEFAULT NULL,
    p_nombre_puesto VARCHAR DEFAULT NULL
)
RETURNS SETOF RECORD AS $$
BEGIN
  IF opcion = 1 THEN
    RETURN QUERY SELECT estado_coleccion
    FROM tarjeta t JOIN equipo e ON t.id_equipo=e.id_equipo
    WHERE e.nombre_equipo = p_nombre_equipo;

  ELSIF opcion = 2 THEN
    RETURN QUERY SELECT medio_adquisicion
    FROM tarjeta t JOIN jugador j ON t.id_jugador=j.id_jugador
    WHERE j.nombre_jugador = p_nombre_jugador;

  ELSIF opcion = 3 THEN
    RETURN QUERY SELECT d.nombre_deporte,e.nombre_equipo,t.fecha_salida_venta
    FROM deporte d JOIN equipo e ON d.id_deporte=e.id_deporte
    JOIN tarjeta t ON e.id_equipo=t.id_equipo
    WHERE d.nombre_deporte = p_nombre_deporte;

  ELSIF opcion = 4 THEN
    RETURN QUERY SELECT nombre_equipo
    FROM equipo WHERE id_equipo IN (
      SELECT id_equipo FROM tarjeta
      WHERE estado_coleccion = p_estado_coleccion);

  ELSIF opcion = 5 THEN
    RETURN QUERY SELECT j.nombre_jugador,p.nombre_puesto
    FROM jugador j JOIN equipo e ON j.id_equipo=e.id_equipo
    JOIN puesto p ON j.id_puesto=p.id_puesto
    WHERE e.nombre_equipo=p_nombre_equipo
      AND p.nombre_puesto=p_nombre_puesto;
  END IF;
END;
$$ LANGUAGE plpgsql;
\end{lstlisting}
\end{frame}


\end{document}
